\documentclass[10pt]{mypackage}

\usepackage[uprightgreeks,noDcommand]{kpfonts}
\renewcommand{\coloneq}{\coloneqq}
\renewcommand{\eqcolon}{\eqqcolon}
%\usepackage{dsfont}
%\renewcommand{\mathbb}{\mathds}
%\usepackage{homework}
\usepackage{notes}

\usepackage[ backend=bibtex, style = alphabetic, sorting=ynt ]{biblatex}
\addbibresource{../all_references.bib}

\usepackage{parskip}
\frenchspacing

\fancyhf{}
\fancyhead[R]{Avinash Iyer}
\fancyhead[L]{Kernels and GNS for Group Representations}
\fancyfoot[C]{\thepage}

\setcounter{secnumdepth}{0}

\begin{document}
\RaggedRight
\section{Unitary Representations}%
\subsection{Group Actions on Measure Spaces}%
Let $G$ be a topological group. The action of $G$ on a measure space $\left( \Omega,\mathcal{M},\mu \right)$ is called measurable if the map $G\times\Omega\to \Omega$ taking $\left( g,\omega \right)\mapsto g.\omega$ is measurable when $G$ is equipped with the Borel $\sigma$-algebra.

We say the measure $\mu$ is invariant if $g_{\ast}\mu(A) = \mu(A)$ for all $A\in \mathcal{M}$. The measure is said to be quasi-invariant if $\mu$ and $g_{\ast}\mu$ are of the same measure class (or mutually absolutely continuous).

If $G\curvearrowright \Omega$ is a measurable group action on a $\sigma$-finite space, or $G$ acts continuously on the locally compact space $\Omega$ with $\mu$ a regular Borel measure, then there exists a positive measurable function, the Radon--Nikodym derivative $\frac{dg_{\ast}\mu}{d\mu}$, on $\Omega$, such that
\begin{align*}
  \int_{}^{} f\left( g.\omega \right)\frac{dg_{\ast}\mu}{d\mu}(\omega)\:d\mu(\omega) &= \int_{}^{} f(\omega)\:d\mu(\omega).
\end{align*}
The function $c_{\mu}\colon G\times \Omega\to [0,\infty]$ given by $c_{\mu}\left( g,\omega \right) = \frac{dg_{\ast}\mu}{d\mu}(\omega)$ is measurable and satisfies the cocycle relation
\begin{align*}
  c_{\mu}\left( g_1g_2,\omega \right) &= c_{\mu}\left( g_1,g_2.\omega \right)c_{\mu}\left( g_{2},\omega \right)
\end{align*}
for all $g_1,g_2\in G$ and a.e. $\omega\in \Omega$.

Define for each $g\in G$ an operator $\pi_g$ on $L^2\left( \Omega,\mu \right)$ by
\begin{align*}
  \pi_gf(\omega) &= c_{\mu}\left( g^{-1},\omega \right)^{1/2}f\left( g^{-1}.\omega \right)
\end{align*}
for all $f\in L^2\left( \Omega,\mu \right)$ and $\omega\in \Omega$. Thus, the operator $\pi_g$ is unitary, and the cocycle relation gives that $\pi\colon G\to \mathcal{U}\left( L^2\left( \Omega,\mu \right) \right)$ is a group homomorphism.
\begin{lemma}
  Suppose $G$ is a locally compact group, and $\mathscr{H}$ is a separable Hilbert space. Let $\pi\colon G\to \mathcal{U}\left( \mathscr{H} \right)$ is a homomorphism from $G$ to the unitary group of $ \mathscr{H} $. If the function $g\mapsto \iprod{\pi(g)\xi}{\xi}$ is measurable for all $\xi\in \mathscr{H}$, then $\pi$ is SOT-continuous.
\end{lemma}
\begin{proof}
  Let $\xi$ be a vector in $ \mathscr{H} $. It suffices to show that $g\mapsto \pi(g)\xi$ is continuous at $e$.

  Let $\ve > 0$ and $A = \set{g\in G : \norm{\pi(g)\xi - \xi} < \ve/2}$. Then, $A$ is measurable since
  \begin{align*}
    A &= \set{g\in G : 2\re\left( \iprod{\pi(g)\xi}{\xi} \right) > 2\norm{\xi}^2-\ve^2/4},
  \end{align*}
  and $A = A^{-1}$ with $A^2 = AA^{-1}\subseteq \set{g\in G : \norm{\pi(g)\xi -\xi} < \ve}$.

  The subset $\pi(G)\xi$ is separable since $ \mathscr{H} $ is separable. Thus, there is a sequence $\left( g_n \right)_n$ in $G$ such that $\left( \pi\left( g_n \right)\xi \right)_{n}$ is dense in $\pi(G)\xi$. For any $g\in G$, we have
  \begin{align*}
    \norm{\pi\left( g_n^{-1}g \right)\xi - \xi} &= \norm{\pi\left( g \right)\xi - \pi\left( g_n \right)\xi}\\
                                                &< \ve/2,
  \end{align*}
  so that $G = \bigcup_{n}g_nA$. Since the Haar measure of $G$ is not zero, the Haar measure of $A$ is not zero, so $A^2 = AA^{-1}$ is a neighborhood of $e$.
\end{proof}
\begin{proposition}
  Suppose $G$ is a $\sigma$-compact locally compact group, and $\left( \Omega,\mu \right)$ is a $\sigma$-finite measurable space such that $L^2\left( \Omega,\mu \right)$ is separable, and let $G\curvearrowright \Omega$ be a measure class preserving action. Then, $\pi\colon G\to L^2\left( \Omega,\mu \right)$ defined by $\pi_gf(\omega) = c\left( g^{-1},\omega \right)^{1/2}f\left( g^{-1}.\omega \right)$ is a unitary representation of $G$ on $L^2\left( \Omega,\mu \right)$.
\end{proposition}
\begin{proof}
  Since $g$ is $\sigma$-compact, the matrix coefficient realization
  \begin{align*}
    g\mapsto \iprod{\pi_gf_1}{f_2} = \int_{}^{} c_{\mu}\left( g^{-1},\omega \right)^{1/2} f_1\left( g^{-1}.\omega \right) \overline{f_2(\omega)}\:d\mu(\omega)
  \end{align*}
  is measurable by Fubini's theorem, hence SOT continuous.
\end{proof}
\section{Kernels and Functions of Positive Type}%
Let $X$ be a topological space
\begin{definition}
  A kernel of positive type on $X$ is a continuous function $\Phi\colon X\times X\to \C$ such that for any $n\in \N$, any elements $x_1,\dots,x_n\in X$, and any $c_1,\dots,c_n\in \C$,
  \begin{align*}
    \sum_{i=1}^{n} \sum_{j=1}^{n} c_i \overline{c_j}\Phi\left( x_i,x_j \right) &\geq 0.
  \end{align*}
  Equivalently, the matrix $\left( \Phi\left( x_i,x_j \right) \right)_{i,j}\in \M_n$ is positive.
\end{definition}
\begin{proposition}
  If $\Phi$ is a kernel of positive type on $X$, then for all $x,y\in X$, one has
  \begin{enumerate}[(i)]
    \item $\Phi\left( y,x \right) = \overline{\Phi\left( x,y \right)}$;
    \item $\left\vert \Phi\left( x,y \right) \right\vert^2 \leq \Phi\left( x,x \right)\Phi\left( y,y \right)$ (i.e., Cauchy--Schwarz).
  \end{enumerate}
\end{proposition}
\begin{proof}
  When $n=2$, $x_1 = x$, and $x_2 = y$, one observes that the matrix
  \begin{align*}
    A &= \begin{pmatrix}\Phi\left( x,x \right) & \Phi\left( x,y \right) \\ \Phi\left( y,x \right) & \Phi\left( y,y \right)\end{pmatrix}
  \end{align*}
  is positive, so it is hermitian with positive diagonal values satisfying $\Phi\left( x,x \right)\Phi\left( y,y \right) - \Phi\left( y,x \right)\Phi\left( x,y \right) \geq 0$.
\end{proof}
The canonical example is the following. If $f\colon X\to \mathscr{H}$ is a continuous map into a Hilbert space, then
\begin{align*}
  \Phi\left( x,y \right) &= \iprod{f(x)}{f(y)}
\end{align*}
is a kernel of positive type on $X$.
\begin{theorem}[GNS for Kernels of Positive Type]
  Let $\Phi$ be a kernel of positive type on a topological space $X$. Then, there exists a Hilbert space $\mathscr{H}$ and a continuous map $f\colon X\to \mathscr{H}$ such that
  \begin{enumerate}[(i)]
    \item $\Phi\left( x,y \right) = \iprod{f(x)}{f(y)}$;
    \item $ \overline{\Span}\set{f(x) : x\in X} = \mathscr{H} $.
  \end{enumerate}
  This Hilbert space is unique up to isomorphism.
\end{theorem}
\begin{proof}
  For every $x\in X$, we let $\Phi_x(y) = \Phi\left( x,y \right)$. These are each continuous functions. Let $V = \Span\set{\Phi_x : x\in X}\subseteq C(X)$. Given two elements of $V$
  \begin{align*}
    \varphi &= \sum_{i=1}^{m}a_i\Phi_{x_i}\\
    \psi &= \sum_{j=1}^{n} b_j\Phi_{x_j},
  \end{align*}
  we write
  \begin{align*}
    \iprod{\varphi}{\psi} &= \sum_{i=1}^{m}\sum_{j=1}^{n} a_i \overline{b_j}\Phi\left( x_i,x_j \right)\\
                          &= \sum_{j=1}^{n} \overline{b_j}\varphi\left( x_j \right)\\
                          &= \sum_{i=1}^{m} a_i \overline{\psi\left( x_i \right)}.
  \end{align*}
  The form $ \iprod{\varphi}{\psi} $ is a positive sesquilinear form on $V$. We claim it is definite. Notice that $ \iprod{\varphi}{\Phi_x} = \varphi(x) $, and by Cauchy--Schwarz,
  \begin{align*}
    \left\vert \varphi(x) \right\vert^2 &\leq \Phi\left( x,x \right) \iprod{\varphi}{\varphi}.
  \end{align*}
  Completing $V$ with respect to the induced norm yields the Hilbert space $ \left( \mathscr{H}, \iprod{\cdot}{\cdot} \right) $.

  If $f\colon X\to V$ takes $x\mapsto \Phi_x$, then $ \iprod{f(x)}{f(y)} = \Phi\left( x,y \right) $, and $f$ is continuous since
  \begin{align*}
    \norm{f(x)-f(y)}^2 &= \Phi\left( x,x \right) + \Phi\left( y,y \right) - 2\re\left( \Phi\left( x,y \right) \right)
  \end{align*}
  and $\Phi$ is continuous.

  Now, suppose $\left( \mathscr{K},g \right)$ is another pair satisfying (i) and (ii). Then,
  \begin{align*}
    \norm{\sum_{i=1}^{n}\Phi_{x_i}}^2 &= \sum_{i=1}^{n}\sum_{j=1}^{n} a_i \overline{a_j}\Phi\left( x_i,x_j \right)\\
                                      &= \sum_{i=1}^{n} \sum_{j=1}^{n} a_i \overline{a_j} \iprod{g\left( x_i \right)}{g\left( x_j \right)}\\
                                      &= \norm{\sum_{i=1}^{n}a_ig\left( x_i \right)}^2.
  \end{align*}
  Thus, if we define the map $V\to \mathscr{K}$ by taking $\Phi_{x_i}\mapsto g\left( x_i \right)$ and extending linearly, we obtain an isometry $T\colon \mathscr{H}\to \mathscr{K}$. Since $T(V) = \Span\set{g(x) : x\in X}$, which is dense in $ \mathscr{K} $, we have that $T\circ f = g$.
\end{proof}
We note that, given any kernels of positive type $\Phi$ and $\Psi$, the following are also kernels of positive type:
\begin{enumerate}[(i)]
  \item $ \overline{\Phi} $;
  \item $ t\Phi $ for all $t\geq 0$;
  \item $\Phi + \Psi$;
  \item $\Phi\Psi$;
  \item the pointwise limit of a family of kernels of positive type, given that the limit is continuous.
\end{enumerate}
Thus, the kernels of positive type are a convex cone in the vector space of continuous kernels on $X$ that is closed under conjugation, pointwise product, and pointwise convergence.
\begin{lemma}
  Let $\Phi$ be a kernel of positive type on $X$ such that $\left\vert \Phi\left( x,y \right) \right\vert < 1$ for all $x,y\in X$. Then, $\left( 1-\Phi \right)^{-t}$ is a kernel of positive type for all $t\geq 0$.
\end{lemma}
\begin{proof}
  The series
  \begin{align*}
    1 + \sum_{j=1}^{\infty}\frac{\prod_{k=0}^{j-1}\left( t+k \right)}{j!}\Phi\left( x,y \right)^{j}
  \end{align*}
  converges to $\left( 1-\Phi\left( x,y \right) \right)^{-t}$ for all $x,y\in X$ by the assumption that $\left\vert \Phi\left( x,y \right) \right\vert < 1$ for all $x,y\in X$.
\end{proof}
A sort of dual to the kernels of positive type is a kernel of conditionally negative type. They are related to each other by Schoenberg's Theorem.
\begin{definition}
  A real-valued kernel $\Psi\colon X\times X \to\R$ is said to be of \textit{conditionally negative type} if
  \begin{enumerate}[(i)]
    \item $\Psi\left( x,x \right) = 0$ for all $x\in X$;
    \item $\Psi\left( x,y \right) = \Psi\left( y,x \right)$ for all $x,y\in X$;
    \item for all elements $x_1,\dots,x_n$ and all real numbers $c_1,\dots,c_n$ with $\sum c_i = 0$, one has
      \begin{align*}
        \sum_{i=1}^{n}\sum_{j=1}^{n}c_ic_j\Psi\left( x_i,x_j \right) &\leq 0.
      \end{align*}
  \end{enumerate}
\end{definition}
The canonical example is that of the map on a real Hilbert space
\begin{align*}
  \Psi\colon \mathscr{H}\times \mathscr{H}&\to \R\\
  \left( \xi,\eta \right) &\mapsto \norm{\xi-\eta}^2.
\end{align*}
One observes that for any $\xi_1,\dots,\xi_n\in \mathscr{H}$ and $c_1,\dots,c_n\in \R$ with $\sum c_i = 0$, we have
\begin{align*}
  \sum_{i=1}^{n}\sum_{j=1}^{n}c_ic_j\Psi\left( \xi_i,\xi_j \right) &= \sum_{i=1}^{n}\sum_{j=1}^{n} c_ic_j \norm{\xi_i-\xi_j}^2\\
                                                                   &= -2\norm{\sum_{i=1}^{n}c_i\xi_i}^2.
\end{align*}
So, if $X$ is a topological space and $\mathscr{H}$ is a real Hilbert space, with $f\colon X\to \mathscr{H}$ a continuous mapping, then the kernel $\left( x,y \right)\mapsto \norm{f(x)-f(y)}^2$ is of conditionally negative type. Similar to how kernels of positive type arise, this is universal via a GNS construction.
\begin{theorem}[GNS for Kernels of Conditionally Negative Type]
  Let $\Psi$ be a kernel of conditionally negative type on $X$, and let $x_0\in X$. There exists a real Hilbert space $ \mathscr{H} $ and a continuous map $f\colon X\to \mathscr{H}$ such that
  \begin{enumerate}[(i)]
    \item $\Psi\left( x,y \right) = \norm{f(x)-f(y)}^2$;
    \item $ \overline{\Span}\set{f(x) - f(x_0) : x\in X} = \mathscr{H}$.
  \end{enumerate}
  Moreover, this pair is unique up to affine isometry.
\end{theorem}
\begin{proof}
  Let $V = \ker\left( \epsilon \right)$, where $\epsilon\colon \R[X]\to \R$ is the augmentation map
  \begin{align*}
    \sum_{i=1}^{m}c_i\delta_{x_i} &\mapsto \sum_{i=1}^{m}c_i.
  \end{align*}
  Let $\varphi = \sum_{i=1}^{m}a_i\delta_{x_i}$ and $\psi = \sum_{j=1}^{n}b_j\delta_{x_j}$ in $V$. Then, we let
  \begin{align*}
    \iprod{\varphi}{\psi} &= -\frac{1}{2}\sum_{i=1}^{n}\sum_{j=1}^{n}a_ib_j\Psi\left( x_i,x_j \right).
  \end{align*}
  This defines a positive symmetric form on $V$. Set
  \begin{align*}
    N_{\Psi} &= \set{\varphi\in V : \iprod{\varphi}{\varphi} = 0}.
  \end{align*}
  By the Cauchy--Schwarz inequality, this is a linear subspace of $V$. Define
  \begin{align*}
    \iprod{\left[ \varphi \right]}{\left[ \psi \right]} &= \iprod{\varphi}{\psi}
  \end{align*}
  on $V/N_{\Psi}$. Let $\left( \mathscr{H}, \iprod{\cdot}{\cdot} \right)$ be the Hilbert space completion of $V/N_{\Psi}$. Set $f\colon X\to H$ taking $x\mapsto \left[ \delta_x-\delta_{x_0} \right]$. Then,
  \begin{align*}
    \norm{f(x) - f(y)}^2 &= \{si\left( x,y \right),
  \end{align*}
  since $\Psi\left( x,x \right) = 0$ and $\Psi\left( y,y \right) = 0$ with $\Psi\left( y,x \right) = \Psi\left( x,y \right)$. Moreover, $f$ is continuous since $\Psi$ is continuous, so we obtain existence.

  Let $\left( \mathscr{K},g \right)$ be another such pair. We claim that the map
  \begin{align*}
    \left[ \sum_{i=1}^{n}a_i\delta_{x_i} \right] &\mapsto \sum_{i=1}^{n} a_ig\left( x_i \right)
  \end{align*}
  extends to an affine isometry. Since $\sum_{i=1}^{n}a_i = 0$, it follows that
  \begin{align*}
    \norm{\left[ \sum_{i=1}^{n}a_i\delta_{x_i} \right]}^2 &= \norm{\left[ \sum_{i=1}^{n}a_i\left( \delta_{x_i}-\delta_{x_i} \right) \right]}^2\\
                                                          &= -\frac{1}{2}\sum_{i,j=1}^{n} a_ia_i\left( \Psi\left( x_i,x_j \right)- \Psi\left( x_i,x_0 \right)-\Psi\left( x_j,x_0 \right) \right)\\
                                                          &= -\frac{1}{2}\sum_{i,j=1}^{n} a_ia_j\left( \norm{g\left( x_i \right)-g\left( x_j \right)}^2 - \norm{g\left( x_i \right)-g\left( x_0 \right)}^2 - \norm{g\left( x_j \right)-g\left( x_0 \right)}^2 \right)\\
                                                          &= \sum_{i,j=1}^{n} a_ia_j \iprod{g\left( x_i \right)-g\left( x_0 \right)}{g\left( x_j \right)-g\left( x_0 \right)}\\
                                                          &= \norm{\sum_{i=1}^{n}a_i\left( g\left( x_i \right)-g\left( x_0 \right) \right)}^2\\
                                                          &= \norm{\sum_{i=1}^{n}a_ig\left( x_i \right)}^2.
  \end{align*}
  Moreover, since $\Span\set{g(x)-g\left( x_0 \right) : x\in X}$ is dense in $ \mathscr{K} $, it follows that $T$ is surjective with $T\circ f = g$.
\end{proof}
Analogously to the case of kernels of positive type, there are certain inheritance properties for kernels of conditionally negative type (i.e., kernels of conditionally negative type form a convex cone and are closed in the topology of pointwise convergence on the space of continuous functions). Furthermore, if $\Phi$ is a real valued kernel of positive type, then $\Phi\left( x,x \right)-\Phi\left( x,y \right)$ is a kernel of conditionally negative type.
\begin{lemma}
  Let $X$ be a topological space, and let $\Psi\colon X\times X\to \R$ be a continuous kernel with $\Psi\left( x,x \right) = 0$ and $\Psi\left( y,x \right) = \Psi\left( x,y \right)$ for all $x,y\in X$. For a fixed $x_0\in X$, let $\Phi\colon X\times X\to \R$ be defined by
  \begin{align*}
    \Phi\left( x,y \right) &= \Psi\left( x,x_0 \right) + \Psi\left( y,x_0 \right) - \Psi\left( x,y \right).
  \end{align*}
  Then, $\Psi$ is a kernel of conditionally negative type if and only if $\Phi$ is a kernel of positive type.
\end{lemma}
\begin{proof}
  Suppose $\Phi$ is a kernel of positive type. Then, for any $x_1,\dots,x_n\in X$ and $c_1,\dots,c_n\in \R$ with $\sum_{i=1}^{n}c_i = 0$, we have
  \begin{align*}
    0 &\leq \sum_{i=1}^{n}\sum_{j=1}^{n}c_ic_j\Phi\left( x_i,x_j \right)\\
      &= -\sum_{i=1}^{n}\sum_{j=1}^{n}c_ic_j\Psi\left( x_i,x_j \right),
  \end{align*}
  showing that $\Psi$ is of conditionally negative type.

  Suppose $\Psi$ is of conditionally negative type. Then, by the GNS for kernels of conditionally negative type, there is a real Hilbert space $\mathscr{H}$ and a continuous map $f\colon X\to \mathscr{H}$ such that $\Psi\left( x,y \right) = \norm{f(x)-f(y)}^2$ for all $x,y\in X$. Then, we have
  \begin{align*}
    \Phi\left( x,y \right) &= \Psi\left( x,x_0 \right) + \Psi\left( y,x_0 \right) - \Psi\left( x,y \right)\\
                           &= 2 \iprod{f(x)-f\left( x_0 \right)}{f\left( y \right)-f\left( x_0 \right)}
  \end{align*}
  for all $x,y\in X$. Thus, for all $x_1,\dots,x_n\in X$ and $c_1,\dots,c_n\in \R$, we have
  \begin{align*}
    \sum_{i,j=1}^{n}c_ic_j\Phi\left( x_i,x_j \right) &= 2\norm{\sum_{i=1}^{n}c_i\left( f\left( x_i \right)-f\left( x_0 \right) \right)}^2\\
                                                     &\geq 0.
  \end{align*}
  Thus, since $\Psi$ is symmetric the expression
  \begin{align*}
    \sum_{i,j=1}^{n} c_i \overline{c_j}\Phi\left( x_i,x_j \right) &\geq 0
  \end{align*}
  holds for all $x_1,\dots,x_n\in X$ and all complex $c_1,\dots,c_n$. Thus, $\Phi$ is of positive type.
\end{proof}
\begin{theorem}[Schoenberg's Theorem]
  Let $X$ be a topological space, $\Psi\colon X\times X\to \R$ a continuous kernel on $X$ with $\Psi\left( x,x \right) = 0$ and $\Psi\left( y,x \right) = \Psi\left( x,y \right)$ for all $x,y\in X$. Then, the following are equivalent:
  \begin{enumerate}[(i)]
    \item $\Psi$ is of conditionally negative type;
    \item $\exp\left( -t\Psi \right)$ is of positive type for all $t\geq 0$.
  \end{enumerate}
\end{theorem}
\begin{proof}
  Suppose $\exp\left( -t\Psi \right)$ is of positive type for every $t\geq 0$. Then, the kernel $1-\exp\left( -t\Psi \right)$ is of conditionally negative type, and so is the pointwise limit
  \begin{align*}
    \Psi &= \lim_{t\to 0} \frac{1-\exp\left( -t\Psi \right)}{t}.
  \end{align*}
  This gives the desired result.

  Suppose now that $\Psi$ is of conditionally negative type. It suffices to show that $\exp\left( -\Psi \right)$ is of positive type. Let $x_0\in X$. Then, the kernel
  \begin{align*}
    \Phi\left( x,y \right) &= \Psi\left( x,x_0 \right) + \Psi\left( y,x_0 \right)-\Psi\left( x,y \right)
  \end{align*}
  is of positive type. Notice then that
  \begin{align*}
    \exp\left( -\Psi\left( x,y \right) \right) &= \exp\left( \Phi\left( x,y \right) \right)\exp\left( -\Psi\left( x,x_0 \right) \right)\exp\left( -\Psi\left( y,x_0 \right) \right).
  \end{align*}
  Since $\Phi$ is of positive type, it is also the case for $\Phi\left( x,y \right)^{n}$, hence $\exp\left( \Phi \right)$ is of positive type. Since
  \begin{align*}
    \sum_{i=1}^{n}\sum_{j=1}^{n} c_i \overline{c_j} \exp\left( -\Psi\left( x_i,x_0 \right) \right)\exp\left( -\Psi\left( x_j,x_0 \right) \right) &= \left\vert \sum_{i=1}^{n}c_i \exp\left( -\Psi\left( x_i,x_0 \right) \right) \right\vert^2\\
                                                                                                                                                                &\geq 0
  \end{align*}
  for all $x_1,\dots,x_n\in X$ and complex numbers $c_1,\dots,c_n$. Thus, $\exp\left( -\Psi \right)$ is of positive type since it is the product of kernels of positive type.
\end{proof}
\section{Functions of Positive Type}%
\subsection{The GNS Construction for Representations}%
\begin{definition}
  A function of positive type on a locally compact group $G$ is a continuous function $\varphi\colon G\to \C$ such that the kernel $\left( g_1,g_2 \right)\mapsto \varphi\left( g_2^{-1}g_1 \right)$ is of positive type.
\end{definition}
\begin{proposition}
  Let $\varphi$ be a function of positive type on $G$. Then, for all $g,h\in G$, we have
  \begin{enumerate}[(i)]
    \item $ \overline{\varphi(g)} = \varphi\left( g^{-1} \right) $
    \item $ \left\vert \varphi(g) \right\vert\leq \varphi(e) $
    \item $\left\vert \varphi(g)-\varphi(h) \right\vert^2 \leq 2\varphi(e)\left( \varphi(e)-\re\left( \varphi\left( g^{-1}h \right) \right) \right)$.
  \end{enumerate}
\end{proposition}
\begin{proof}
  Part (i) follows from the fact that 
  \begin{align*}
    \varphi\left( g \right) &= \varphi\left( (g)(e) \right)\\
                            &= \overline{\varphi\left( (e)\left( g^{-1} \right) \right)}\\
                            &= \overline{\varphi\left( g^{-1} \right)},
  \end{align*}
  and part (ii) follows from
  \begin{align*}
    \left\vert \varphi\left( g \right) \right\vert^2 &= \left\vert \varphi\left( (g)(e) \right) \right\vert\\
                                                     &\leq \varphi\left( g^{-1}g \right)\varphi\left( e^{-1}e \right)\\
                                                     &= \varphi\left( e \right)^2.
  \end{align*}
  Finally, (iii) follows from the fact that for all $x,y,z\in \C$, we have
  \begin{align*}
    \left\vert x \right\vert^2 + \left\vert y \right\vert^2 + \left\vert z \right\vert^2 + 2\re \left( \varphi\left( g^{-1} \right) \overline{y}x + \varphi\left( h^{-1} \right) \overline{z}x + \varphi\left( g^{-1}h \right) \overline{y}z \right) &\geq 0
  \end{align*}
  since $\varphi$ is of positive type. When 
  \begin{align*}
    x &= \left\vert \varphi(g)-\varphi(h) \right\vert\\
    y &= - \frac{\left\vert \varphi(g)-\varphi(h) \right\vert}{\varphi(g)-\varphi(h)}\\
    z &= \frac{\left\vert \varphi(g)-\varphi(h) \right\vert}{\varphi(g)-\varphi(h)},
  \end{align*}
  we have
  \begin{align*}
    \left\vert \varphi(g)-\varphi(h) \right\vert^2 + 2 + 2\re\left( \overline{\left( -\varphi(g) + \varphi(h) \right)}\frac{\left\vert \varphi(g)-\varphi(h) \right\vert^2}{ \overline{\varphi(g)-\varphi(h)} }-\varphi\left( g^{-1}h \right) \right) &\geq 0.
  \end{align*}
\end{proof}
\begin{proposition}
  Let $\left( \pi, \mathscr{H} \right)$ be a unitary representation of $G$, and let $\xi$ be a vector in $ \mathscr{H} $. Then, the matrix coefficient $g\mapsto \iprod{\pi(g)\xi}{\xi}$ is a function of positive type.
\end{proposition}
\begin{proof}
  Immediate from the discussed example.
\end{proof}
\begin{definition}
  Let $G$ be a locally compact group with Haar measure $dx$. Given functions $f,g$ on $G$, the convolution $f\ast g$ is given by
  \begin{align*}
    f\ast g(x) &= \int_{}^{} f\left( xy \right)g\left( y^{-1} \right)\:dy\\
               &= \int_{}^{} f\left( y \right)g\left( y^{-1}x \right)\:dy.
  \end{align*}
\end{definition}
\begin{proposition}
  Every matrix coefficient of the regular representation of a locally compact group $G$ vanishes at infinity.
\end{proposition}
\begin{proof}
  We note that the matrix coefficients of the left regular representation of $G$ are given by
  \begin{align*}
    \set{f\ast \check{g} : f,g\in L^2\left( G \right)},
  \end{align*}
  where $\check{g}(x) = g\left( x^{-1} \right)$. Since $C_c(G)\subseteq L^2(G)$ is dense, there exist $f_n,g_n\in C_c(g)$ such that $\lim f_n = f$ and $\lim g_n = g$. So, $\lim f_n\ast \check{g}_n = f\ast \check{g}$ uniformly on $G$.

  Thus, to show that $f\ast\check{g}\in C_0(G)$, we may analogously assume that $f,g\in C_c(G)$. If $K = \supp(f)$ and $L = \supp(g)$, then $\supp\left( f\ast\check{g} \right) \subseteq KL^{-1}$. Thus, $f\ast\check{g}\in C_c(G)$.
\end{proof}
\begin{corollary}
  Let $G$ be a non-compact locally compact group. Then, $\lambda_G$ has no finite dimensional subrepresentation.
\end{corollary}
\begin{proof}
  Let $\left( \pi,\mathscr{H} \right)$ be a finite-dimensional unitary representation of $G$, and let $\left( v_i \right)_i$ be an orthonormal basis. It suffices to show that some matrix coefficient $ \iprod{\pi\left( \cdot \right)v_i}{v_j} $  does not vanish at infinity.

  Yet, observe that $\left( \iprod{\pi(g)v_i}{v_j} \right)_{i,j}$ is the matrix of $\pi(g)$ with respect to the basis $\left( v_i \right)_i$, so it is a unitary matrix. Thus, for any $i$, we have
  \begin{align*}
    \sum_{j} \left\vert \iprod{\pi(g)v_i}{v_j} \right\vert^2 &= 1
  \end{align*}
  for all $g\in G$.
\end{proof}

\begin{definition}
  A unitary representation $\left( \pi,\mathscr{H} \right)$ of $G$ is cyclic if there exists $\xi\in \mathscr{H}$ such that $ \overline{\Span}\set{\pi_g\xi : g\in G} $ is dense in $ \mathscr{H} $. We call $\xi$ a cyclic vector for $\pi$.
\end{definition}
\begin{theorem}[GNS for Topological Groups]
  Let $\varphi$ be a function of positive type on a topological group $G$. Then, there exists a triple $\left( \pi,\mathscr{H},\xi \right)$ consisting of a cyclic unitary representation with cyclic vector $\xi$ such that
  \begin{align*}
    \varphi(g) &= \iprod{\pi(g)\xi}{\xi}.
  \end{align*}
  This triple is unique in the sense that if $\left( \rho,\mathscr{K},\eta \right)$ is another such triple, then there exists an intertwining Hilbert space isomorphism $T\colon \mathscr{H}\to \mathscr{K}$ such that $T\xi = \eta$.
\end{theorem}
\begin{proof}
  Let $\Phi$ be a kernel of positive type associated to $\varphi$. Then, there exists a Hilbert space $ \mathscr{H} $ and a continuous map $f\colon G\to \mathscr{H}$ such that
  \begin{align*}
    \iprod{f(x)}{f(y)} &= \Phi\left( x,y \right)\\
                       &= \varphi\left( y^{-1}x \right).
  \end{align*}
  We note that for every $g\in G$, $ \iprod{f(gx)}{f(gy)} = \iprod{f(x)}{f(y)} $. Thus, by uniqueness, there is a unitary operator $\pi_g$ on $ \mathscr{H} $ such that $\pi_gf(x) = f(gx)$ for all $x\in G$. By density of $\set{f(x) : x\in G}$, we have that $\pi_{g_1g_2} = \pi_{g_1}\pi_{g_2}$ for all $g_1,g_2\in G$. Since $f$ is continuous, $g\mapsto \pi_g\eta$ is continuous for every $\eta\in \mathscr{H}$. Set $\xi\coloneq f(e)$. Then, $\xi$ is cyclic with $\varphi(g) = \iprod{\pi_g\xi}{\xi}$ for all $g\in G$.

  If $\left( \rho,\mathscr{H},\eta \right)$ is a triple with the same properties, then we have that
  \begin{align*}
    \Phi\left( g_1,g_2 \right) &= \iprod{\pi_{g_1}\eta}{\pi_{g_2}\eta},
  \end{align*}
  By uniqueness, there exists an isomorphism $T\colon \mathscr{H}\to \mathscr{K}$ such that $T\pi_g\xi = \pi_g\eta$ for all $g\in g$. Thus, $T\xi = \eta$, and $T$ intertwines $\pi$ and $\rho$ since $\xi$ is cyclic.
\end{proof}
\begin{example}
  If $G$ is a discrete group, and $\varphi = \delta_e$ is the Dirac mass at $e$, then the GNS representation for $\varphi$ is the regular representation for $G$.
\end{example}
\subsection{The Cone of Functions of Positive Type}%
We let $\mathcal{P}(G)$ be the set of functions of positive type on the topological group $G$. This forms a convex cone closed under complex conjugation and pointwise product. We let $ \mathcal{P}_1(G) $ be the subset of $\varphi\in \mathcal{P}(G)$ such that $\varphi(e) = 1$.
\begin{proposition}
  Let $\varphi_1,\varphi_2\in \mathcal{P}(G)$, and let $\varphi = \varphi_1 + \varphi_2$. Then, $\pi_{\varphi_1}\leq \pi_{\varphi}$.

  Moreover, if $\pi_{\varphi}$ is irreducible, then $\varphi_1 = t\varphi$ for some $t\geq 0$.
\end{proposition}
\nocite{kazhdan_property_t,conway_functional_analysis}
\printbibliography
\end{document}
